%
% Vorlage für Versuchsprotokolle im Grundpraktikum an der RWTH Aachen
% -------------------------------------------------------------------
%
% Bitte verwenden Sie diese Vorlage zur Erstellung Ihrer
% Protokolle.
%
% Wenn möglich, drucken Sie Ihre Protokolle bitte doppelseitig
% aus. Wenn Sie das Protokoll doppelseitig ausdrucken, verwenden Sie
% die twoside Option ...
\documentclass[twoside]{protokoll}
\usepackage{float}
% ... wenn Sie nicht doppelseitig ausdrucken können, lassen Sie sie
% weg:
%\documentclass{protokoll}
 
% Setzen Sie hier den jeweiligen Versuchsnamen ein:
\versuchsname{Versuchsname}

% Setzen Sie hier Ihre(n) Namen und Ihre Matrikelnummer(n) sowie Ihre
% Gruppennummer ein:
\teilnehmer{Sophie, 466344}
\teilnehmer{Alexa, 469147}
\gruppe{C4}

\begin{document}

% Hier wird die zum Versuch gehörende Aufgaben ausgewählt, indem das
% Prozentzeichen (%) vor dem entsprechenden \input-Befehl entfernt
% wird. Die tex-Datei kann im Moodle beim jeweiligen Versuch
% heruntergeladen werden. Achten Sie darauf, dass sich alle benötigten
% Dateien im selben Verzeichnis befinden, bevor Sie diese tex-Datei
% kompilieren.
%
% Die Bearbeitung der Aufgaben findet unmittelbar in der jeweiligen
% Datei ("aufgaben_xyz.tex") statt. Dort befinden sich Umgebungen
% der Form
%
% \begin{aufgabe}
%   ... (Aufgabentext) ...
% \end{aufgabe}
%
% Hinter dem versuchsspezifischen Aufgabentext und dem "\end{aufgabe}"
% beginnen Sie jeweils Ihre Lösung der Aufgabe.
%


% Thermodynamik
\section{Adiabatenkoeffizient der Luft}

\begin{aufgabe}{Analyse der Kugelschwingungen}
  Beschreiben Sie kurz Ihr Vorgehen bei der Messung. Zeigen Sie
  exemplarisch für jeden Erlenmeyerkolben sowie die große Mariottesche
  Flasche je einen Schwingungsverlauf und seine Auswertung. Tabellieren
  Sie Ihre aus allen Schwingungsverläufen ermittelten Ergebnisse und
  die zugehörigen Kugelpositionen im Rohr.
\end{aufgabe}

% Bitte belassen Sie die Aufgabentexte in Ihrem Protokoll und beginnen
% Sie hier mit der Lösung der ersten Aufgabe:

% Achten Sie darauf, alle verwendeten Formeln anzugeben und die
% verwendeten Formelzeichen zu definieren.
\section{Overview and rudemintary measurements}

We first fix the flask in place and attach the tubes tightly to make sure no air escapes. we then connect the glass tube to the top of the flask and make sure it is vertical. we make sure the pressure sensor is connected and the whole setup is ready to conduct the experiment, we then lower the ball to a position where fall damage is not a concern using the magnets. once we reach a suitable spot, we start conducting the experiment by releasing the ball and recording the pressure change until the ball reaches the equilibrium position, we then measure the hight of the ball using the ruler and repeat the experiment for few times.

\vspace{1em}
we summarise the measurements in the following tables:
\begin{table}[h]
  \begin{tabular}{l|ccc}
    Properties & \multicolumn{3}{|c}{Measurements } \\ \\ \hline
     m (g) & 16,5 & 16,4 & 16,4  \\
           &  &  &  \\ \hline
     p (hPa)& 981 & 980 &  \\
           &  &  &  \\ \hline
     D (mm) & 15,98 &15,55  &  \\
           &  &  &  \\ \hline
     T$_{0}$ (s) &0,38  & 0,39 & 0,38 \\
           & 0,39 &  0,38 &0,39  \\ \hline
  \end{tabular}
  \caption{Measurements for the 1 Liter flask.}
  \label{tab:beispieltabelle}
\end{table}

\begin{itemize}
  \item We measured the mass of the ball using the scale available in the lab.
  \item We measured the pressure using the pressure sensor in lab along the first few minutes.
  \item We measured the diameter of the ball using the micrometer caliper available in the lab.
  \item We measured the period of the oscillation using the Cassy measurements and the time difference between multiple peaks.
\end{itemize}

\section{Errors and assumptions}

\begin{table}[h]
  \centering
  \begin{tabular}{l|ccc}
    Measurement tool & Stat. error  & Sys. error & Combined errors\\ \\ \hline
     Scale (g) &  &  &  \\
           &  &  &  \\ \hline
     Pressure sensor (hPa)&  &  &  \\
           &  &  &  \\ \hline
     Micrometer caliper (mm) &  &  &  \\
           &  &  &  \\ \hline
     Cassy (s) &  &  &  \\
           &  &  &  \\ \hline
  \end{tabular}
  \caption{Beispieltabelle.}
  \label{tab:beispieltabelle}
\end{table}







\begin{aufgabe}{Messungen der übrigen Observablen}
  Beschreiben Sie die Messungen der übrigen Observablen. Geben Sie
  jeweils Ihre Mess\-er\-geb\-nis\-se an und bestimmen Sie die
  zugehörigen Messunsicherheiten. Verwenden Sie ggfs. Tabellen, um
  Ihre Ergebnisse übersichtlich darzustellen.
\end{aufgabe}









\begin{aufgabe}{Ergebnis und Fehlerrechnung}
  Berechnen Sie $\kappa$ für die einzelnen Messungen. Beschreiben
  Sie, wie Sie die Unsicherheit auf den Adiabatenkoeffizienten
  $\kappa$ bestimmen. Tabellieren Sie exemplarisch für die große
  Mariottesche Flasche die Größe der einzelnen Beiträge zur
  Unsicherheit auf $\kappa$. Geben Sie dabei jeweils an, ob eine
  Unsicherheit von Messung zu Messung unkorreliert ist, oder ob es
  sich um eine korrelierte Unsicherheit handelt. Geben Sie bei den
  korrelierten Unsicherheiten an, ob sie nur für ein einzelnes
  Gasvolumen oder über alle Messungen hinweg korreliert sind.
  
  Stellen Sie die für die einzelnen Messungen erhaltenen Werte von
  $\kappa$ in einem Plot dar. Machen Sie die von Messung zu Messung
  unkorrelierten Fehler jeweils durch einen Fehlerbalken kenntlich und
  die (geeignet zusammengefassten) korrelierten Fehler durch
  Bänder. Geben Sie Ihr Endergebnis für $\kappa$ an. Vergleichen Sie
  mit dem Literaturwert und diskutieren Sie Ihr Ergebnis.
\end{aufgabe}


\begin{table}[h]
  \centering
  \begin{tabular}{l|ccc}
    Properties & \multicolumn{3}{|c}{Measurements } \\ \\ \hline
     M (g) &  &  &  \\
           &  &  &  \\ \hline
     p(hPa)&  &  &  \\
           &  &  &  \\ \hline
     D(mm) &  &  &  \\
           &  &  &  \\ \hline
  \end{tabular}
  \caption{Beispieltabelle.}
  \label{tab:beispieltabelle}
\end{table}



%
% Der nachfolgende Abschnitt mag nützlich sein, um Vorlagen für häufig
% vorkommende Arbeitsschritte mit LaTeX daraus zu entnehmen. Er muss
% vor Abgabe des Protokolls entfernt werden. (Das \end{document} muss
% aber natürlich erhalten bleiben.)
%
%
\end{document}
